\section{ODE and state-space methods}
\label{app:ode}

The bridge between integral equations and differential equations: when the kernel
is built from exponentials, the Volterra equation \emph{is} an ODE in disguise.

\subsection{The exponential kernel gives a first-order ODE}

Take the convolution MBPP equation $\xi=s+\phi\conv\xi$ with
$\phi(t)=\kappa\theta e^{-\theta t}$, and split off the endogenous part
$y:=\xi-s=\phi\conv\xi$. Differentiating
$y(t)=\kappa\theta\int_0^t e^{-\theta(t-u)}\xi(u)\dd u$ (Leibniz rule),
\begin{equation}
y'(t)=\kappa\theta\,\xi(t)-\theta\,y(t)=(\kappa-1)\theta\,y(t)+\kappa\theta\,s(t),
\qquad y(0)=0,\quad \xi=s+y .
\label{eq:exp-ode}
\end{equation}
So $\xi$ obeys a \emph{first-order linear ODE forced by the baseline $s$}, with a
\emph{stable} pole at $(\kappa-1)\theta<0$. Solving \eqref{eq:exp-ode} with the
integrating factor reproduces the single-exponential resolvent
$h(t)=\kappa\theta e^{(\kappa-1)\theta t}$.

\begin{intuition}
A single memory timescale is a single pole, i.e.\ a first-order filter, i.e.\ a
first-order ODE. One number $\kappa$ sets the feedback gain and one number $\theta$
sets the relaxation rate $(1-\kappa)\theta$: push the system and it returns to
baseline at that rate.
\end{intuition}

\subsection{Rational Laplace transform $\Leftrightarrow$ finite ODE}

\begin{theorem}[Realisation]\label{thm:rational}
The convolution Volterra equation $\xi=s+\phi\conv\xi$ is equivalent to a
finite-order linear constant-coefficient ODE (forced by $s$ and its derivatives)
\emph{if and only if} $\phi$ has a strictly proper rational Laplace transform ---
equivalently, iff $\phi$ is a finite exponential-polynomial
$\sum_q p_q(t)e^{-b_q t}$.
\end{theorem}
\begin{proof}
($\Leftarrow$) If $\Lap\{\phi\}=P/Q$ with $\deg P<\deg Q=:d$, then from
\eqref{eq:transfer} $\Lap\{\xi\}=\frac{Q}{Q-P}\Lap\{s\}$, so
$(Q-P)(z)\Lap\{\xi\}=Q(z)\Lap\{s\}$. As $Q-P$ has degree $d$, reading
$z^k\leftrightarrow d^k/\dd t^k$ yields the order-$d$ ODE
$(Q-P)(\tfrac{d}{dt})\xi=Q(\tfrac{d}{dt})s$. ($\Rightarrow$) Conversely a finite ODE
makes the transfer function $Q/(Q-P)$ rational, hence
$\Lap\{\phi\}=1-(Q-P)/Q=P/Q$ rational; and a function has rational Laplace transform
iff it is an exponential-polynomial (partial fractions and inverse Laplace).
\end{proof}

\begin{intuition}
A memory built from finitely many decaying modes can be tracked by finitely many
state variables --- ``rational transfer function $\Leftrightarrow$
finite-dimensional machine'' (this is realisation theory). A power-law memory is
\emph{not} rational: it carries infinitely many timescales, so no finite ODE exists
(one then uses Volterra quadrature, Appendix~\ref{app:volterra}).
\end{intuition}

\subsection{Sums of exponentials: the state-space realisation}

For $\phi(\tau)=\sum_q a_q e^{-b_q\tau}$ the natural states are
$y_q=\bigl(a_q e^{-b_q\cdot}\bigr)\conv\xi$, giving the constant-coefficient system
\begin{equation}
y_q'=a_q\,\xi-b_q\,y_q,\qquad \xi=s+\sum_q y_q,
\end{equation}
i.e.\ $\dot u=Au+a\,s$, $\xi=s+\mathbf 1^\top u$ with $A=a\mathbf 1^\top-\diag(b)$.
The multivariate version (kernel matrix) carries an index pair $(m,j)$ and gives the
$M^2$-state system \eqref{eq:statespace}, implemented as
\code{solve\_mbpp\_ode\_multivariate}. Since sums of exponentials are dense, this
realises (or approximates) essentially any kernel.

\subsection{Numerical integration and the time-varying case}

These ODEs are integrated by a standard explicit Runge--Kutta (RK4) sweep on the
observation grid; the compensator $\Xi=\int_0^t\xi$ is accumulated alongside (an
extra trapezoidal state). The constant-coefficient systems are autonomous in their
coefficients; the \emph{covariate-modulated} system \eqref{eq:ltv} is identical in
form but with the readout gain $\alpha^0_{mj}e^{\delta_{mj}^\top Z(t)}$ depending on
$t$ --- a \emph{linear time-varying} (LTV) ODE.

\begin{proposition}[LTV well-posedness]
If the coefficient maps $t\mapsto A(t)$ and $t\mapsto a(t)$ are continuous (here
through the continuous covariate path $Z(t)$), the LTV system $\dot u=A(t)u+a(t)s(t)$
has a unique solution on $[0,T]$ for every initial condition, given by the
time-ordered (Peano--Baker) series; numerically, RK4 converges at its usual order.
\end{proposition}

\begin{intuition}
A time-varying linear ODE is no harder to \emph{integrate} than a constant one ---
you just re-read the coefficients at each step. What you lose is the \emph{closed
form}: there is no single matrix exponential $e^{At}$ when $A$ depends on $t$ (the
$A(t)$ at different times need not commute), which is the differential-equations
echo of ``no constant transfer function'' from Appendix~\ref{app:transforms}. So
\code{solve\_mbpp\_ltv} integrates \eqref{eq:ltv} step by step, and that is exactly
how the excitation-covariate model of \S\ref{sec:excitation-cov} is computed.
\end{intuition}
